\documentclass[12pt]{article}
\usepackage[utf8]{inputenc}
\usepackage[spanish]{babel}
\usepackage{amsmath, amssymb}
\usepackage{geometry}
\geometry{a4paper, margin=2.5cm}
\usepackage{xcolor}
\usepackage{tcolorbox}
\tcbuselibrary{skins}

\definecolor{azulIPn}{HTML}{003DA5}
\definecolor{verdeMatem}{HTML}{1de9b6}
\definecolor{rojoSuper}{HTML}{ff6b6b}
\definecolor{moradoVec}{HTML}{9b59b6}

\newtcolorbox{definicion}[1]{
  colback=azulIPn!10, colframe=azulIPn,
  fonttitle=\bfseries, title=Definición: #1
}
\newtcolorbox{ejemplo}{
  colback=verdeMatem!10, colframe=verdeMatem!60!black,
  fonttitle=\bfseries, title=Ejemplo
}
\newtcolorbox{nota}{
  colback=rojoSuper!10, colframe=rojoSuper!60!black,
  fonttitle=\bfseries, title=Nota importante
}

\title{\textcolor{azulIPn}{\textbf{Nota 03 — Superficies Paramétricas}} \\
       \large Cálculo Vectorial aplicado a animaciones}
\author{Daniel Tovar Abraham \\ ESCOM -- IPN}
\date{Servicio Social 2026}

\begin{document}
\maketitle

% -----------------------------------------------
\section{¿Qué es una superficie paramétrica?}
% -----------------------------------------------

\begin{definicion}{Superficie paramétrica}
  Una superficie paramétrica es una función vectorial 
  $\vec{r}: D \subset \mathbb{R}^2 \to \mathbb{R}^3$ que mapea 
  dos parámetros $(u, v)$ a un punto en el espacio:
  \[
    \vec{r}(u,v) = \bigl(x(u,v),\; y(u,v),\; z(u,v)\bigr)
  \]
\end{definicion}

% -----------------------------------------------
\section{Vectores tangentes y normal}
% -----------------------------------------------

Para una superficie paramétrica $\vec{r}(u,v)$ se definen:

\subsection{Vectores tangentes parciales}
\[
  \vec{r}_u = \frac{\partial \vec{r}}{\partial u} = 
  \left(\frac{\partial x}{\partial u},\;
        \frac{\partial y}{\partial u},\;
        \frac{\partial z}{\partial u}\right)
\]
\[
  \vec{r}_v = \frac{\partial \vec{r}}{\partial v} = 
  \left(\frac{\partial x}{\partial v},\;
        \frac{\partial y}{\partial v},\;
        \frac{\partial z}{\partial v}\right)
\]

\subsection{Vector normal}
\[
  \hat{n} = \frac{\vec{r}_u \times \vec{r}_v}{\|\vec{r}_u \times \vec{r}_v\|}
\]

\begin{nota}
  Si en algún punto $\vec{r}_u \times \vec{r}_v = \vec{0}$, 
  la superficie tiene un \textbf{punto singular} — no se puede 
  definir la normal ahí.
\end{nota}

% -----------------------------------------------
\section{Ejemplo: Cinta de Möbius}
% -----------------------------------------------

\begin{ejemplo}
La Cinta de Möbius se parametriza como:
\[
  \vec{r}(u,v) = \left(
    \left(R + v\cos\tfrac{u}{2}\right)\cos u,\;
    \left(R + v\cos\tfrac{u}{2}\right)\sin u,\;
    v\sin\tfrac{u}{2}
  \right)
\]
con $u \in [0, 2\pi]$, $v \in [-1, 1]$ y $R = 2$.

\medskip
\textbf{¿Por qué es especial?}
\begin{itemize}
  \item Tiene \textcolor{rojoSuper}{\textbf{solo una cara}} — si recorres 
        toda la superficie regresas al punto de inicio pero 
        en el lado opuesto
  \item Tiene \textcolor{rojoSuper}{\textbf{solo un borde}} — un único 
        lazo continuo
  \item Es \textcolor{azulIPn}{\textbf{no orientable}} — el vector normal 
        $\hat{n}$ no se puede definir de forma continua en toda la superficie
\end{itemize}
\end{ejemplo}

% -----------------------------------------------
\section{Ejemplos de otras superficies}
% -----------------------------------------------

\subsection{Esfera}
\[
  \vec{r}(\phi, \theta) = 
  \bigl(R\sin\phi\cos\theta,\; R\sin\phi\sin\theta,\; R\cos\phi\bigr)
\]
con $\phi \in [0,\pi]$, $\theta \in [0, 2\pi]$.

\subsection{Toro}
\[
  \vec{r}(u,v) = 
  \bigl((R + r\cos v)\cos u,\; (R + r\cos v)\sin u,\; r\sin v\bigr)
\]
con $u, v \in [0, 2\pi]$, $R$ = radio mayor, $r$ = radio menor.

\subsection{Paraboloide}
\[
  \vec{r}(u,v) = \bigl(u,\; v,\; u^2 + v^2\bigr)
\]

\bigskip
\begin{nota}
  Todas estas superficies se pueden animar en Manim usando \texttt{Surface()} 
  con la misma parametrización que aparece aquí.
\end{nota}

\end{document}